The canonical form only includes ``lesser or equal''
inequalities. Therefore, we can transform each equality constraint
into two inequality constraints. 


The first constraint is
\begin{equation}
  \label{eq:c1}
  x_1 + x_2 + x_3 = 1.
\end{equation}
It can be written with two inequality constraints:
\begin{align*}
  x_1 + x_2 + x_3 &\leq 1, \\
  x_1 + x_2 + x_3 &\geq 1, \\
\end{align*}
or, equivalently, 
\begin{align*}
  x_1 + x_2 + x_3 &\leq 1, \\
  -x_1 - x_2 - x_3 &\leq -1. \\
\end{align*}
The second constraint is
\begin{equation}
  \label{eq:c2}
x_1 - x_2 + x_4 = 1.
\end{equation}
It can be written with two inequality constraints:
\begin{align*}
x_1 - x_2 + x_4 &\leq 1, \\
x_1 - x_2 + x_4 &\geq 1, \\
\end{align*}
or, equivalently, 
\begin{align*}
x_1 - x_2 + x_4 &\leq 1, \\
-x_1 + x_2 - x_4 &\leq -1. \\
\end{align*}

Finally, the fact that  $x_i$ is non negative can be
written as $-x_i \leq 0$, for $i=1,\ldots,4$.

Putting everything together, we obtain
\begin{align*}
  x_1 + x_2 + x_3 &\leq 1, \\
  -x_1 - x_2 - x_3 &\leq -1, \\
x_1 - x_2 + x_4 &\leq 1, \\
-x_1 + x_2 - x_4 &\leq -1, \\
-x_1 & \leq 0,\\
-x_2 & \leq 0,\\
-x_3 & \leq 0,\\
-x_4 & \leq 0.
\end{align*}
In matrix form, we have
\[
 \mathcal{P} = \{x \in \R^4 \mid Ax \leq b\},
\]
where
\[
A = \begin{pmatrix*}[r]
  1 & 1 & 1 & 0 \\
  -1 & -1 & -1 & 0  \\
  1 & -1 & 0 & 1 \\
  -1 & 1 & 0 & -1 \\
  -1 & 0 & 0 & 0  \\
  0 & -1 & 0 & 0  \\
  0 & 0  & -1 & 0  \\
  0 & 0  &  0 & -1  \\
\end{pmatrix*}, \; b = \rvect{1 \\ -1 \\ 1 \\ -1 \\ 0 \\ 0 \\ 0 \\ 0  }
\]

\clearpage
Note that there are other ways to write the polyhedron in canonical
form. In particular, if we observe that $x_3$ and $x_4$ in the
standard form are slack variables, we can eliminate them in the
canonical form. Indeed, we can combine \req{eq:c1} with the non negativity constraint of $x_3$ to obtain:
\[
x_3 = 1 - x_1 - x_2 \geq 0,
\]
or, equivalently
\[
x_1 + x_2 \leq 1.
\]
And $x_3$ disappears from the formulation. 
Similarly, \req{eq:c2} can be written
\[
x_4 = 1 - x_1 + x_2 \geq 0,
\]
or, equivalently
\[
x_1 - x_2 \leq 1.
\]
Putting everything together, we obtain:
    \begin{align*}
        x_1 + x_2 \leq 1, \\
        x_1 - x_2 \leq 1, \\
        - x_1 \leq 0, \\
        - x_2 \leq 0. \\
    \end{align*}

In matrix form, it is:
    \[
    \mathcal{P} = \{x \in \R^2 \mid Ax \leq b\}
    \]
    with
    \[
    A = \begin{pmatrix*}[r] 1 & 1 \\  1 & -1 \\  -1 & 0 \\  0 & -1 \end{pmatrix*}, \qquad b = \begin{pmatrix*}[r] 1 \\ 1 \\ 0 \\ 0 \end{pmatrix*}. 
    \]
This is now a polyhedron in $\R^2$.
